\section{Programme}\label{sec:prog}

\subsection*{Track A: lowering the bound}
\begin{enumerate}
\item \emph{Climb the ladder.} The certified pipeline now stands at $k=50$
($H_1\le246$, the established bound, via the $\varepsilon$-enlarged simplex).
Stadlmann's region ($k=49$, $240$): version~0.6 reproduces the variational
step exactly (Theorem~\ref{thm:capped-exact}, Corollary~\ref{cor:240}), so the
project's bound is $240$ \emph{conditionally on the unverified
equidistribution estimates of \cite{Sta26}}. For $212$ ($k=45$,
\cite{Cha26}) version~0.7 confirms the variational criterion for the caps of
their Table~3 (Proposition~\ref{prop:cha-capped}); before the project's bound
can move, the discrepancy between their Table~3 and Lemma~9.1 must be
resolved; version~0.8 has the tool for optimising on the capped support
(Proposition~\ref{prop:adjoint}), and for the caps of Lemma~9.1 no polynomial in
our $846$-dimensional space reaches $3.99$ (Proposition~\ref{prop:cha-ceiling});
so the draft's value must belong to its Table~3 caps. For $236$ ($k=48$) the variational criterion of \cite{Kin26} holds with
our own degree-$22$ vector (Proposition~\ref{prop:kin-capped}), and the
capped numerator of the published witness of \cite{SY26} is reproduced
exactly, integer for integer, by two independent programs
(Proposition~\ref{prop:sy-exact}), so the project's bound is $236$
conditionally (Corollary~\ref{cor:236}); the analytic input remains
\lbl{Claimed}. Our own vector is confirmed by
two programs (version~0.12); Proposition~\ref{prop:bonf} makes such checks cheaper
when only the inequality is needed. Then
$186$ and $182$.
\item \emph{Floor of the method (our main original direction).} For the plain
simplex the dual bound $M_k\le\frac{k}{k-1}\log k$ shows where the method
stops ($k\ge51$). For the $\varepsilon$-enlarged simplex
Theorem~\ref{thm:floor} gives $k\ge30$, which is still far from the expected
truth near $50$; Section~4 explains what a sharper weight must do, and
Proposition~\ref{prop:family} indicates that it must depend on more of $t$
than $t_m$ and $S$ (a Collatz--Wielandt weight $\psi_m=\int F\,dt_m/F$ built
from a good positive test function $F$ is the natural candidate). For the
one-band hybrid supports of \cite{Sta26,Cha26},
Theorem~\ref{thm:hybrid-floor} gives the first such bound we know of
($k\ge29$ at Stadlmann's level, $k\ge27$ at that of \cite{Cha26}), still
far from $49$ and $45$; it ignores the caps. A rigorous
dual certificate there would say whether $182$ is close to the limit of the
present inputs or far from it, and which input (which Type~I/II/III range) is
binding.
\item \emph{Which rung is cheap.} From Table~\ref{tab:tuples}, the gain per
unit of $k$ below the claimed frontier is $182\to176\to168\to162\to158$
for $k=39,\dots,35$: the steps $39\to38$ and $38\to37$ are worth the most.
\end{enumerate}

\begin{problem}[\lbl{Open}; partly answered by Theorem~\ref{thm:hybrid-floor} for one band, without using the caps]
Give a rigorous upper bound for the variational quantity of \cite{Sta26} as a
function of $k$ and of the admissible exponents, and deduce the least $k$ the
hybrid method can reach with the currently proved estimates.
\end{problem}

\subsection*{Track B: the parity barrier}
No sieve of the above kind can distinguish integers with an even number of
prime factors from those with an odd number. Known ways around the barrier
all inject bilinear or structural information. We will study, one per
session, what each would need in order to say something about twin primes:
Chen's theorem \cite{Che73}; Bombieri's asymptotic sieve \cite{Bom76}; the
parity-breaking sieve of Friedlander and Iwaniec \cite{FI98}; Heath-Brown's
theorem that Siegel zeros would imply the conjecture \cite{HB83}; the
function-field proof of Sawin and Shusterman \cite{SS22}; and the
entropy-based results on correlations of multiplicative functions
\cite{Tao16}.

\begin{problem}[\lbl{Open}, to be made precise in Lean]
State, as a single explicit hypothesis on bilinear sums over $n(n+2)$, the
weakest known input that implies the twin prime conjecture, and formalise the
implication.
\end{problem}

\subsection*{Rules of this paper}
A numerical bound enters the ``established by this project'' column only with
an exact certificate and, for the combinatorial part, a Lean proof. Ideas are
logged with their outcome, including failures. Nothing here is peer reviewed.

